\chapter{Guía rápida para editar la plantilla}
\label{cap:anexo-guia}

Este anexo resume los cambios habituales al comenzar un TFG con la plantilla.

\section{Datos del documento}
\label{sec:datos-documento}

Los datos de portada y los metadatos del PDF se completan al principio de
\texttt{main.tex}:

\begin{itemize}
    \item \texttt{\textbackslash titulo\{Título en español\}}: título del TFG en español.
    \item \texttt{\textbackslash tituloIngles\{Title in English\}}: título en inglés.
    \item \texttt{\textbackslash autor\{Nombre y apellidos\}}: nombre y apellidos.
    \item \texttt{\textbackslash titulacion\{Grado cursado\}}: grado cursado.
    \item \texttt{\textbackslash director\{Nombre y apellidos\}}: director o directora.
    \item \texttt{\textbackslash codirector\{Nombre y apellidos\}}: codirección; se deja vacío si no existe.
    \item \texttt{\textbackslash fechaEntrega\{9\}\{2026\}}: mes numérico y año de entrega.
\end{itemize}

La universidad, el centro, el tipo de documento y la ciudad ya están
configurados para la EIMIA. El título completo se reutiliza en la portada, el
resumen y los metadatos.

\section{Organización de archivos}
\label{sec:organizacion-archivos}

\begin{description}
    \item[\texttt{main.tex}] Datos del TFG y orden de las partes del documento.
    \item[\texttt{packagesTFG.tex}] Formato y comandos internos de la plantilla.
    \item[\texttt{preambulo/}] Portada, créditos, dedicatoria, índices, agradecimientos, resumen y abstract.
    \item[\texttt{caps/}] Un archivo por capítulo o anexo.
    \item[\texttt{figs/}] Imágenes y logotipos.
    \item[\texttt{bibliography.bib}] Referencias bibliográficas en formato BibTeX.
\end{description}

Para añadir un capítulo, crea su archivo en \texttt{caps/} y añade el
\texttt{\textbackslash include} correspondiente en \texttt{main.tex}.
La dedicatoria, los agradecimientos y el bloque de anexos pueden desactivarse
comentando las líneas señaladas en ese archivo.

\section{Compilación}
\label{sec:compilacion}

La plantilla utiliza pdfLaTeX y BibTeX. Puedes compilarla desde la terminal,
en Overleaf o con un editor local.

\begin{lstlisting}[language=bash,caption={Compilación automática},label={lst:compilacion}]
latexmk -pdf -interaction=nonstopmode main.tex
\end{lstlisting}

Este comando es la opción recomendada en la terminal porque ejecuta las pasadas
necesarias. Si no dispones de \texttt{latexmk}, usa \texttt{pdflatex},
\texttt{bibtex} y dos pasadas adicionales de \texttt{pdflatex}.

\begin{description}
    \item[Overleaf.] Sube el proyecto, selecciona \texttt{main.tex} como
    documento principal y usa pdfLaTeX. No requiere una instalación local.
    \item[Texmaker o TeXstudio.] Abre \texttt{main.tex}, establécelo como
    documento maestro y configura la compilación rápida con \texttt{latexmk}
    o con la secuencia manual anterior.
    \item[Visual Studio Code.] Abre la carpeta del proyecto y utiliza una
    receta de LaTeX Workshop basada en \texttt{latexmk}.
\end{description}

Si aparecen referencias \texttt{??}, vuelve a compilar. Si persiste un error,
consulta el primer mensaje relevante de \texttt{main.log}.

\section{Lista de comprobación antes de entregar}
\label{sec:checklist-entrega}

\begin{itemize}
    \item No quedan textos de ejemplo ni datos provisionales.
    \item La portada contiene título español e inglés, autoría, titulación, dirección y fecha correctos.
    \item El resumen y el abstract tienen entre 250 y 350 palabras y entre 3 y 5 palabras clave.
    \item Todas las figuras y tablas están citadas y señalan su fuente cuando procede.
    \item Todas las referencias citadas aparecen en la bibliografía y siguen el estilo requerido.
    \item No quedan referencias cruzadas sin resolver.
    \item Las páginas del cuerpo y de las referencias están numeradas en el pie.
    \item Se han revisado ortografía, estilo, pies de figura y formato de tablas.
    \item Se han preparado por separado los documentos adicionales exigidos por la convocatoria.
\end{itemize}
